\item \textbf{Demostración de un sistema de lentes zoom.}
Un sistema de lentes de zoom mantiene posiciones fijas del objeto y la pantalla de proyección mientras varía la magnificación al desplazar una o más lentes intermedias. Monte en un banco óptico un objeto, dos lentes convergentes y una pantalla. La primera lente tiene una distancia focal $f_1 = 5.00\text{ cm}$ y la segunda tiene $f_2 = 10.0\text{ cm}$. La pantalla está a la derecha de la segunda. Inicialmente, el objeto se sitúa a $7.50\text{ cm}$ a la izquierda de la primera lente, y la imagen final proyectada en la pantalla tiene un aumento de $+1.00$.
\begin{enumerate}
    \item Determine la distancia total entre el objeto y la pantalla.
    \item Ahora se desplazan ambas lentes a lo largo del eje óptico de manera que la imagen sigue enfocada en la pantalla pero la magnificación lateral es $+3.00$. Determine el desplazamiento neto de cada lente desde su posición inicial.
    \item ¿Existe más de una combinación posible de desplazamientos para lograr esto?
\end{enumerate}